\documentclass[11pt]{amsart}
\usepackage[a4paper,textwidth=158mm,textheight=237mm,centering]{geometry}
\usepackage[T1]{fontenc}
\usepackage{lmodern,microtype,mathtools,amssymb,booktabs}
\usepackage[unicode,colorlinks=true,linkcolor=black,citecolor=black,urlcolor=blue]{hyperref}
\newtheorem{theorem}{Theorem}
\newtheorem{lemma}{Lemma}
\newcommand{\Rpos}{\mathbb R_{>0}}
\newcommand{\one}{\mathbf 1}
\newcommand{\diag}{\operatorname{diag}}
\title[Sharp equilibrium counts with cubic sources]{Sharp equilibrium counts with cubic sources}
\author{Harvey Yorke}\address{Valyu AI}
\ifdefined\journalversion
\hypersetup{pdfauthor={Harvey Yorke}}
\else
\author{Valyu DeepResearch}\address{Valyu AI}
\hypersetup{pdfauthor={Harvey Yorke and Valyu DeepResearch}}
\fi
\hypersetup{pdftitle={Sharp equilibrium counts with cubic sources}}
\date{18 October 2026}
\newcommand{\paperAuthorBlock}[1]{%
  \begin{minipage}[t]{.44\textwidth}\centering\normalfont\normalsize
  #1\par\small Valyu AI\par\end{minipage}}
\makeatletter
\renewcommand{\@setauthors}{%
  \begingroup\trivlist\centering\normalfont\normalsize
  \@topsep30\p@\relax\advance\@topsep by -\baselineskip\item\relax
  \ifdefined\journalversion\paperAuthorBlock{Harvey Yorke}%
  \else\paperAuthorBlock{Harvey Yorke}\hspace{.06\textwidth}%
  \paperAuthorBlock{Valyu DeepResearch}\fi\endtrivlist\endgroup}
\renewcommand{\@setaddresses}{}
\makeatother
\begin{document}
\raggedbottom
\begin{abstract}
A mass action network with $n$ species, $n$ reactions and stoichiometric
rank $n-1$ has at most $n$ positive nondegenerate equilibria on a
stoichiometric class. We construct a family attaining this bound for every
$n\ge2$, with positive rational rate constants and maximum source
molecularity three. The construction answers the sharpness question of
Banaji and Feliu; their quadratic obstruction then proves that source
molecularity three is minimal. The equilibrium set is an explicit
monomial curve. Sparse polynomial interpolation gives exactly $n$ simple
intersections with one class, and transversality proves nondegeneracy.
The count persists under small changes of rates and class level.
The classes are unbounded, product molecularities may exceed three,
and simultaneous stability is not asserted.
\end{abstract}
\maketitle
\begin{center}\small
Research draft. Exact checks and internal proof reviews are available;
author approval, priority assessment and external review remain open.
\end{center}

\section{The sharpness question}
An $(n,m,r)$ mass action network has $n$ species, $m$ reactions and
stoichiometric rank $r$. Banaji and Feliu prove that every $(n,n,n-1)$
network has at most $n$ positive nondegenerate equilibria on a
stoichiometric class \cite[Corollary~4.21(i)]{banajiFeliu2026}.
Their Remark~4.23 asks whether the bound is achieved for general $n$
and, if so, what minimum source molecularity is needed.
We answer both questions with a family whose sources have molecularity
at most three. Attainment for $n=2$ is already established in their
Remark~4.23; the contribution is a uniform construction for all $n\ge3$
and the resulting sharp source-molecularity threshold.

The monomial parametrisation of equilibria and the upper bound are part
of the existing theory. The construction below combines a fixed cubic
source matrix, sparse interpolation, and a choice of valid reaction
products to attain the bound. The source degree stays three as $n$ grows,
although the degree of the scalar polynomial used in the construction
does not stay bounded.

An anonymous, unrefereed manuscript \cite{anonymous2026} presents a quadratic
$(n,n+2,n)$ family attaining $n+1$ positive nondegenerate equilibria,
addressing the full-rank question in Banaji and Feliu's Remark~5.8.
The present construction addresses their separate one-conservation-law
question in Remark~4.23 and its minimum source molecularity.

Write reaction $i$ as $y_i\to y_i+\gamma_i$, with nonnegative integral
source and product vectors. Let $A$ have rows $y_i^\top$ and $\Gamma$
have columns $\gamma_i$. The mass action vector field is
\begin{equation}\label{eq:field}
 \dot x=\Gamma v(x),\qquad v_i(x)=\kappa_i x^{y_i},\qquad x\in\Rpos^n.
\end{equation}
Its positive stoichiometric class through $x$ is
$(x+\operatorname{im}\Gamma)\cap\Rpos^n$. An equilibrium is
nondegenerate when the Jacobian restricted to
$\operatorname{im}\Gamma$ is invertible. The full Jacobian is singular
when a conservation law is present.

\begin{theorem}\label{thm:main}
For every integer $n\ge2$, there exists an $(n,n,n-1)$ mass action
network with positive rational rate constants and exactly $n$ positive
equilibria on one stoichiometric class, all nondegenerate. Its maximum
source molecularity is three, the least possible for attaining this
count. Its product molecularities are at most $2n$. The class is unbounded
and its conservation vector has both signs.
\end{theorem}

\section{Cubic sources and a sparse interpolation polynomial}
Let $e_1,\ldots,e_n$ be the standard basis, and fix the sources
\begin{equation}\label{eq:sources}
 y_i=2e_i+e_{i+1}\quad(1\le i<n),\qquad y_n=2e_n.
\end{equation}
The source matrix $A$ is upper bidiagonal with diagonal entries two and
superdiagonal entries one, so $\det A=2^n$. Put
\begin{equation}\label{eq:powers}
 M=2^n,\qquad
 p_i=\frac{2^n+(-1)^{n-i}2^{i-1}}3\quad(1\le i\le n).
\end{equation}
These are distinct positive integers: the two terms in the numerator
cancel modulo three, its value is positive, and the signed powers
$(-1)^{n-i}2^{i-1}$ are distinct. They satisfy
\begin{equation}\label{eq:weighted}
 2p_i+p_{i+1}=M\ (i<n),\qquad 2p_n=M,
 \qquad Ap=M\one.
\end{equation}

Fix distinct positive rational numbers $r_1<\cdots<r_n$; taking $r_j=j$
gives a concrete family for every $n$. Form the rational matrix
\begin{equation}\label{eq:interpolation}
 B_{j0}=1,\qquad B_{ji}=r_j^{p_i}\quad(1\le i\le n).
\end{equation}
Choose a primitive integer vector $\alpha=(\alpha_0,\ldots,\alpha_n)$
in its kernel, with the coefficient of the largest exponent positive.
The next lemma shows that this choice is unique. It can be computed
explicitly from the signed maximal minors of $B$, followed by division
by their greatest common divisor.

\begin{lemma}\label{lem:interpolation}
The polynomial
\begin{equation}\label{eq:polynomial}
 P(t)=\alpha_0+\sum_{i=1}^n\alpha_i t^{p_i}
\end{equation}
has all $n+1$ coefficients nonzero. Its positive roots are exactly
$r_1,\ldots,r_n$, each simple. The coefficients alternate in sign when
ordered by increasing exponent.
\end{lemma}
\begin{proof}
Any $n$ columns of $B$ are independent. Otherwise their dependence
would define a nonzero polynomial with at most $n$ monomials and with
$n$ distinct positive roots, contrary to Descartes' rule of signs.
Thus $B$ has rank $n$ and its one-dimensional kernel contains no vector
with a zero coordinate. The polynomial \eqref{eq:polynomial} has
$n+1$ nonzero monomials and already has $n$ distinct positive roots.
Descartes' rule bounds its positive roots, counted with multiplicity,
by $n$. There are therefore no further positive roots, each prescribed
root is simple, and all $n$ possible coefficient sign changes occur.
\end{proof}

Define
\begin{equation}\label{eq:signs}
 c_i=|\alpha_i|>0,\qquad z_i=\operatorname{sign}(\alpha_i),
 \qquad \rho_i=z_i z_{i+1}\in\{-1,1\}\quad(i<n).
\end{equation}
The $n$ nonconstant coefficients alternate in their exponent order,
so $z$ has both signs for $n\ge2$.

\section{Valid products and the complete equilibrium set}
Choose the stoichiometric reaction vectors
\begin{equation}\label{eq:reactionvectors}
 \gamma_i=-e_i+\rho_i e_{i+1}\quad(i<n),\qquad
 \gamma_n=-\sum_{i=1}^{n-1}\gamma_i.
\end{equation}
For $i<n$ the product is $e_i+(1+\rho_i)e_{i+1}$, a nonnegative
integral vector of molecularity one or three. The final product is
\begin{equation}\label{eq:lastproduct}
 y_n+\gamma_n=e_1+\sum_{j=2}^{n-1}(1-\rho_{j-1})e_j
                   +(2-\rho_{n-1})e_n.
\end{equation}
It is also nonnegative and integral, with molecularity at most $2n$.
All source and product complexes are therefore valid, and the maximum
source molecularity is three. The sources are pairwise distinct and
every $\gamma_i$ is nonzero: its $i$th entry is $-1$ for $i<n$, while
the first entry of $\gamma_n$ is $1$. Thus there are exactly $n$
distinct, nontrivial reactions.

The first $n-1$ columns of $\Gamma$ are independent, as their first
$n-1$ rows form a triangular matrix with diagonal $-1$.
By construction $\Gamma\one=0$ and $z^\top\Gamma=0$. Hence
\begin{equation}\label{eq:stoichiometry}
 \operatorname{rank}\Gamma=n-1,\qquad
 \ker\Gamma=\mathbb R\one,\qquad
 \operatorname{im}\Gamma=\ker z^\top.
\end{equation}
Choose the positive rational rate constants
\begin{equation}\label{eq:rates}
 \kappa_i=c^{-y_i}=\prod_{j=1}^n c_j^{-A_{ij}}.
\end{equation}
At a positive equilibrium, \eqref{eq:stoichiometry} implies
$v(x)=\lambda\one$ for some $\lambda>0$. Taking componentwise
logarithms gives
\[
 A\log(x/c)=(\log\lambda)\one.
\]
Since $A^{-1}\one=p/M$, this is equivalent to
\begin{equation}\label{eq:curve}
 x_i(t)=c_i t^{p_i}\quad(1\le i\le n),\qquad t=\lambda^{1/M}>0.
\end{equation}
Conversely, every point on \eqref{eq:curve} has all reaction fluxes
equal to $t^M$ and is an equilibrium. This is the full positive
equilibrium set, not merely a collection of selected solutions.

Take the nonempty positive stoichiometric class
\begin{equation}\label{eq:class}
 \mathcal S=\{x\in\Rpos^n:z^\top x=-\alpha_0\}.
\end{equation}
It contains $x(r_1)$, and \eqref{eq:stoichiometry} shows that it is
exactly one class. Along the equilibrium curve,
\begin{equation}\label{eq:intersection}
 z^\top x(t)+\alpha_0=P(t).
\end{equation}
Lemma~\ref{lem:interpolation} therefore gives precisely the $n$
positive equilibria $x(r_1),\ldots,x(r_n)$ on $\mathcal S$.

\section{Nondegeneracy and minimal source molecularity}
At $x(t)$ the Jacobian of \eqref{eq:field} is
\begin{equation}\label{eq:jacobian}
 J(t)=t^M\Gamma A\diag\bigl(x_1(t)^{-1},\ldots,x_n(t)^{-1}\bigr).
\end{equation}
The factors to the right of $\Gamma$ are invertible, so this matrix has
rank $n-1$. Differentiating $\Gamma v(x(t))=0$ gives $J(t)x'(t)=0$.
Since $x'(t)\ne0$, its kernel is exactly $\mathbb R x'(t)$.
At each prescribed root,
\begin{equation}\label{eq:transverse}
 z^\top x'(r_j)=P'(r_j)\ne0.
\end{equation}
Thus the kernel of $J(r_j)$ meets $\ker z^\top$ only at zero.
The image of $J(r_j)$ lies in that same $(n-1)$-dimensional space,
so its restriction is invertible. All $n$ equilibria are nondegenerate.

Because $z$ has both signs, there exists $w>0$ with $z^\top w=0$.
For example, give each positive-sign coordinate the number of negative
signs, and each negative-sign coordinate the number of positive signs.
Then $x(r_1)+sw\in\mathcal S$ for all $s\ge0$, proving unboundedness.

The construction has maximum source molecularity three. Conversely,
\cite[Corollary~5.16]{banajiFeliu2026} proves that quadratic
$(n,n,n-1)$ networks have at most two positive nondegenerate equilibria
on a class, which rules out the count $n$ for $n\ge3$. The same
corollary explicitly states that the smallest $n$ for which two are
attained by a quadratic network is three, ruling out attainment for
$n=2$ as well. Therefore no network in the specified class with
maximum source molecularity at most two attains $n$. This proves
the minimality assertion and completes Theorem~\ref{thm:main}.

\smallskip\noindent\emph{Persistence under parameter perturbations.}
The count is not confined to the chosen rational rate vector. Keep the
reaction network fixed and let $\kappa>0$ vary. Its full positive
equilibrium curve is still $x_i(t)=c_i(\kappa)t^{p_i}$, where
\[
 c(\kappa)=\exp\bigl(-A^{-1}\log\kappa\bigr).
\]
On the class $z^\top x=K$, the scalar equation is
$\sum_i z_i c_i(\kappa)t^{p_i}-K=0$.
The implicit function theorem preserves the $n$ simple positive roots
for all $(\kappa,K)$ sufficiently close to the constructed pair.
This polynomial has at most $n+1$ monomials, so Descartes' rule excludes
any additional positive root, including a degenerate one. Thus exactly
$n$ positive nondegenerate equilibria persist on an open neighbourhood
in rate-constant and class-level space. No quantitative neighbourhood
size or dynamical stability assertion is needed.

\section{A three-species example and the scope of the result}
For $n=3$, the powers are $(p_1,p_2,p_3)=(3,2,4)$. Choosing roots
$1,2,3$ yields
\begin{equation}\label{eq:examplepoly}
 P(t)=11t^4-60t^3+85t^2-36
     =(t-1)(t-2)(t-3)(11t+6).
\end{equation}
Thus $c=(60,85,11)$ and $z=(-1,1,1)$. The reactions and rates are
\begin{align}\label{eq:example}
 2X+Y&\xrightarrow{\ 1/306000\ }X,\nonumber\\
 2Y+Z&\xrightarrow{\ 1/79475\ }Y+2Z,\\
 2Z&\xrightarrow{\ 1/121\ }X+2Y+Z.\nonumber
\end{align}
The full positive equilibrium curve is
$(X,Y,Z)=(60t^3,85t^2,11t^4)$, with common reaction flux $t^8$.
The class $-X+Y+Z=36$ has exactly the three equilibria below.
\begin{center}
\begin{tabular}{@{}ccccc@{}}
\toprule
$t$ & $X$ & $Y$ & $Z$ & $P'(t)$\\
\midrule
1 & 60 & 85 & 11 & 34\\
2 & 480 & 340 & 176 & $-28$\\
3 & 1620 & 765 & 891 & 78\\
\bottomrule
\end{tabular}
\end{center}
The final column certifies transversality and nondegeneracy, not stability.
The last product has molecularity four, illustrating why a bound on source
molecularity is different from a bound on every reaction complex.

The conserved linear quantity has mixed signs, so it is not a conserved
positive total mass. Our family is consistent with the upper bound two
for bounded classes in \cite[Corollary~4.21(ii)]{banajiFeliu2026};
for $n>2$, unboundedness is necessary to attain $n$. The result counts
equilibria and does not assert $n$ stable states, a closed mass-balanced
chemical mechanism, or an experimental implementation.

As separate finite checks, run from this directory:
\begin{quote}\small\ttfamily
python3 -B verification/verify\_cubic\_sharpness.py\par
python3 -B -O verification/independent\_cubic\_check.py
\end{quote}
The first reconstructs the sources, products, rational rates and
restricted Jacobians for $2\le n\le8$, checking all 35 prescribed
equilibria with exact integer and rational arithmetic. The second
independently differentiates the mass action equations using only
standard-library rational arithmetic. Results are recorded in
\texttt{verification/cubic-sharpness-certificate.json}. These checks
corroborate the formulas; the all-$n$ theorem rests on the proof above.

\ifdefined\journalversion
\section*{Research assistance}
Valyu DeepResearch supported literature retrieval and source comparison.
Coding agents developed the construction and analytic
argument, implemented the exact verification programs, and drafted and
revised the manuscript. The supplied finite checks corroborate the formulas;
they do not replace the all-$n$ proof or constitute external peer review.
\fi

\bibliographystyle{amsplain}
\bibliography{references}
\end{document}
